\section{Kaon Gluon PDFs from the Pseudo-PDF Method}\label{sec:results}

The Ioffe-time pseudo-distribution (pITD)~\cite{Radyushkin:2017cyf,Orginos:2017kos} is given by
%
\begin{equation}
\mathcal{M}(\nu,z^2) = \langle 0 (P_z)|{\cal O}(z)|0 (P_z)\rangle,
\label{eq:ME_unpol}
\end{equation}
%
The reduced pITD (RpITD)~\cite{Orginos:2017kos,Zhang:2018diq,Li:2018tpe} was constructed to remove the ultraviolet divergences in the pITD by taking the ratio of the pITD to its corresponding $z$-dependent matrix element with $P_z=0$, then normalizing the ratio by the matrix element at $z^2=0$,
 %
\begin{equation}
\mathscr{M}(\nu,z^2)=\frac{\mathcal{M}(zP_z,z^2)/\mathcal{M}(0\cdot P_z,0)}{\mathcal{M}(z\cdot 0,z^2)/\mathcal{M}(0\cdot 0,0)}.
\label{eq:RITD}
\end{equation}
%
The renormalization of ${\cal O}(z)$ and kinematic factors are cancelled in the RpITDs.
By construction, the RpITD double ratios employed here are normalized to one at $z=0$.

Figure~\ref{fig:ensemble_momentum_comparison} shows the lattice-spacing dependence of the kaon gluon RpITD.
The RpITDs are plotted as a function of Wilson link length $z$ at approximately the same momentum $P_z$ for ensembles a12m310 and a15m310, which have lattice spacings around 0.12 and 0.15~fm.
We observe that the RpITD from the a12m310 ensemble has higher values for each  Wilson-link length and has a slower rate in decrease of its values.
Most points are within one sigma of their corresponding points on the other ensemble, with slightly larger central values for most $z$.
As the Wilson-link length increases, we see the difference increases.


\begin{figure}[!tbp]
\centering
\includegraphics[width=0.45\textwidth]{images/ensemble_momentum_comparison.pdf}
\caption{\label{fig:ensemble_momentum_comparison}
Kaon RpITDs for the a12m310 and a15m310 ensembles at boost momentum $P_z \approx 1.3$~GeV as a function of the Wilson-line displacement $z$.
}
\end{figure}


Using the matrix elements obtained from Eq.~\ref{eq:3ptC} and Eq.~\ref{eq:RITD}, we compute the RpITD for both ensembles at the various Wilson link lengths $z$ and meson momenta $P_z$.
Figure~\ref{fig:RITD_fitbands} shows these results compared to the RpITD previously calculated for the pion on the a12m310 ensemble in Ref.~\cite{Fan:2021bcr}.
We see that the kaon RpITDs for both ensembles, shown in the left and middle plots, are congruent.
Additionally, the RpITD for the kaon and pion, shown in the rightmost plot, from a12m310 are very similar, in accordance with the results from Ref.~\cite{Roberts:2021nhw}, where the ratio of the kaon PDF to the pion PDF was calculated and found to be around 1.


\begin{figure*}[!tbp]
\centering
\includegraphics[width=0.32\textwidth]{images/Kaon_RITD_a15m310.pdf}\hfill
\includegraphics[width=0.32\textwidth]{images/Kaon_RITD_a12m310.pdf}\hfill
\includegraphics[width=0.32\textwidth]{images/Pion_RITD_a12m310.pdf}
\caption{\label{fig:RITD_fitbands}
Kaon RpITD for the a15m310 (left), a12m310 (middle) ensembles, and pion RpITD for a12m310 (right).
The bands are the gluon PDF fits to each $z$ by minimizing the $\chi^2$ defined in Eq.~\ref{eq:chi2-RITD-fit} to obtain the gluon PDFs.}
\end{figure*}


We can then extract the gluon PDF distribution through the pseudo-PDF matching condition~\cite{Balitsky:2019krf} that connects the RpITD $\mathscr{M}$ to the lightcone gluon PDF $g(x,\mu^2)$ through
%
\begin{align}
\mathscr{M}(\nu,z^2)&=\int_0^1 dx \frac{xg(x,\mu^2)}{\langle x \rangle_g}R_{gg}(x\nu,z^2\mu^2),
\label{eq:matching-gg}
\end{align}
%
where $\mu$ is the renormalization scale in the $\overline{\text{MS}}$ scheme
and $\langle x \rangle_g=\int_0^1 dx \, x g(x,\mu^2)$ is the gluon momentum fraction of the kaon.
$R_{gg}$ is the gluon-in-gluon matching kernel
%
\begin{equation}
R_{gg}(y,z^2,\mu^2)
=\cos y -\frac{\alpha_s(\mu)}{2\pi}N_c
\left\{ \left[\ln\left(z^2\mu^2\frac{e^{2\gamma_E+1}}{4}\right)+2\right]R_B(y)+R_L(y)+R_C(y)\right\}, \nonumber
\label{eq:Rgg}
\end{equation}
%
where $\alpha_s$ is the strong coupling at scale $\mu$,
$N_c=3$ is the number of colors,
and $\gamma_E=0.5772$ is the Euler-Mascheroni constant.
For the term $R_{gg}(y,z^2,\mu^2)$, $z$ was chosen to be $2e^{-\gamma_E-1/2}/\mu$ so that the logarithmic term vanishes, which suppresses the residuals containing higher order logarithmic terms, following the previous publication regarding the one-loop evolution of the pseudo-PDF~\cite{Radyushkin:2018cvn}.
Equation~\ref{eq:matching-gg} and the terms $R_B(y)$, $R_L(y)$, $R_C(y)$ are defined in Eqs.~7.21--23 and in the paragraph below Eq.~7.23 in Ref.~\cite{Balitsky:2019krf}.

Note that the lattice-calculated RpITDs are also connected to the singlet quark-PDF $q_s$ of the kaon via the quark-gluon matching kernel $R_{gq}$ with an additional $\frac{P_z}{P_0}\int_0^1 dx \frac{xq_S(x,\mu^2)}{\langle x \rangle_g}R_{gq}(x\nu,z^2\mu^2)$ term added to Eq.~\ref{eq:matching-gg}.
In the past pion gluon PDF study~\cite{Fan:2021bcr}, it was found that the quark PDF contribution is small;
therefore in this work, we will neglect the kaon quark PDF.
One can obtain the gluon PDF $g(x,\mu^2)$ by fitting the RpITD through the matching condition in Eq.~\ref{eq:matching-gg};
a similar procedure has also been used by HadStruc Collaboration~\cite{Joo:2019bzr,Joo:2020spy,HadStruc:2021wmh}.

To obtain the gluon PDF $g(x,\mu^2)$ on the right-hand side of Eq.~\ref{eq:matching-gg}, we adopt the phenomenologically motivated form
%
\begin{align}
f_g(x,\mu) = \frac{xg(x, \mu)}{\langle x \rangle_g(\mu)} = \frac{x^A(1-x)^C}{B(A+1,C+1)},
\label{functional}
\end{align}
%
for $x\in[0,1]$ and zero elsewhere.
The beta function $B(A+1,C+1)=\int_0^1 dx\, x^A(1-x)^C$ is used to normalize the area to unity.
Such a form is also used in global fits to obtain the pion gluon PDF by JAM~\cite{Barry:2018ort,Cao:2021aci}.
%
We fit the lattice RpITDs $\mathscr{M}^\text{lat} (\nu,z^2,a,M_\pi)$ obtained in Eq.~\ref{eq:RITD} to the parametrization form $\mathscr{M}^\text{fit}(\nu,\mu,z^2,a,M_\pi)$ in Eq.~\ref{eq:matching-gg} by minimizing the $\chi^2$ function,
%%
\begin{equation}
\label{eq:chi2-RITD-fit}
    \begin{split}
    & \chi^2(\mu,a,M_\pi)=\\
    & \sum _{\nu,z} \frac{(\mathscr{M}^{\rm fit}(\nu,\mu,z^2,a,M_\pi)-\mathscr{M}^{\rm lat}(\nu,z^2,a,M_\pi))^2}{\sigma^2_{\mathscr{M}}(\nu,z^2,a,M_\pi)}.
    \end{split}
\end{equation}
The reconstructed fit bands of the kaon RpITDs for the a15m310 and a12m310 ensembles, and pion RpITDs for a12m310 at each $z^2$, compared with the lattice calculation points, are shown in Fig.~\ref{fig:RITD_fitbands} from left to right.
We see almost no $z^2$-dependence (labeled in different colors) in the reconstructed bands in both a12m310 ensembles, but slightly more dependence in the a15m310 case.
The RpITD data points fluctuate more in the a15m310 kaon data and as a result the fit quality is not as good as in the a12m310 case.
We found the a12m310 fit to both the pion and kaon gluon PDF to have very stable quality with $\chi^2/\text{dof}$ around 1 with consistent output of $f_g(x,\mu)$, regardless of the choice of the maximum value of the Wilson-line displacement $z$.
However, for the a15m310 ensemble, $\chi^2/\text{dof}$ can go as large as 9(3).
By reducing the fitted $z$ to smaller number $z=1$, $\chi^2/\text{dof}$ can be reduced to 5(3).
We suspect that higher-twist effects are enhanced at this coarse lattice spacing such that the fit fails to accurately describe the lattice data.
Possible future work including NNLO matching may help to improve the fit on this ensemble.



\begin{figure*}[!tbp]
  \centering
\includegraphics[width=0.45\textwidth]{images/KaonPDF_Lattice_DSE.pdf}\hfill
\includegraphics[width=0.45\textwidth]{images/xg-x0-1_pion_kaon_a12m310_comparison.pdf}
\caption{\label{fig:DSE_comparison}
%
Left plot: The kaon gluon PDF $xg(x, \mu)/\langle x \rangle_g$ as a function of $x$ obtained from the fit to the lattice data on ensembles with lattice spacing $a\approx\{0.12,0.15\}$~fm (inset plot), pion masses $M_\pi\approx310$~MeV at $a\approx 0.12$~fm, compared with the kaon gluon PDF from DSE'20 at $\mu=2$~GeV in the $\overline{\text{MS}}$ scheme.
%
Right plot: Comparison of pion and kaon gluon PDF $xg(x, \mu)/\langle x \rangle_g$ as a function of $x$ with lattice spacing $a\approx0.12$~fm (outer) and 0.15~fm (inset), pion masses $M_\pi\approx310$~MeV, at $\mu=2$~GeV in the $\overline{\text{MS}}$ scheme.
}
\label{fig:xgx}
\end{figure*}


Our results on the kaon and pion gluon PDFs $xg(x, \mu)/\langle x \rangle_g$ as a function of $x$ from fitting the lattice data are shown in Fig.~\ref{fig:xgx}.
In the inset of the left-hand plot, we show the kaon $xg(x, \mu)/\langle x \rangle_g$ at $\mu=2$~GeV in the $\overline{\text{MS}}$ scheme with $M_\pi\approx 310$~MeV from two lattice spacings $a\approx\{0.12,0.15\}$~fm.
Note that we use $z=1$ for the coarser lattice spacing due to the fit quality becoming significantly worse as $z$ increases, whereas we are able to get a good fit for the finer-lattice RpITD data for $z$ up to 5.
With these choices of data input, our best determination of the kaon $xg(x, \mu)/\langle x \rangle_g$ are consistent within the statistical errors.
We also compare our finer-lattice result with the same quantity obtained from the Dyson-Schwinger equation (DSE) approach ~\cite{Cui:2020tdf}, depicted as the purple band on the left-hand side.
%
Reference~\cite{Cui:2020tdf} provides recent DSE calculations of the pion gluon PDF and the ratio of kaon to pion gluon PDFs as function of $x$.
Although the kaon gluon PDF is not directly provided, it can be obtained from the product of the pion gluon PDF and the kaon-to-pion ratio.
Accordingly, we can make a direct comparison between our lattice calculations and the kaon gluon PDF from DSE calculations.
On the other hand, the DSE'20 error shown is likely overestimated due to lack of correlations in the published parameters for their kaon and pion PDFs.
%
Our kaon gluon PDF results from the smallest lattice spacing, $a\approx 0.12$~fm, is consistent with the DSE results~\cite{Cui:2020tdf} within one-sigma error for $x>0.15$.
On the right-hand side of Fig.~\ref{fig:xgx}, we compare the obtained kaon result at smaller lattice spacing with the pion results obtained from the same ensembles;
we note the kaon gluon PDF is slightly smaller than the one obtained for the pion, similar to its corresponding quark valence PDF and DSE~\cite{Cui:2020tdf} results.
%
We predict the kaon $\langle x^2 \rangle_g/\langle x \rangle_g$ and  $\langle x^3 \rangle_g/\langle x \rangle_g$ to be 0.0779(94) and 0.0187(42), while in good agreement with corresponding results from DSE~\cite{Cui:2020tdf}: 0.075 and 0.015.
On the pion $\langle x^2 \rangle_g/\langle x \rangle_g$, our 310-MeV results gives
 .092(15) and 0.0250(75),
while results from DSE~\cite{Cui:2020tdf},  JAM~\cite{Barry:2018ort,Cao:2021aci} and xFitter~\cite{Novikov:2020snp}
are
0.076, 0.103, 0.158
and 0.015, 0.024, 0.048,
respectively.
%
Future study including finer lattice spacing and lighter pion mass will be important to refine this calculation and provide better predictions on this poorly known meson quantity.
